\documentclass{beamer}
\usepackage{tikz}
\usetheme{Madrid}
\usepackage{algorithm}
\usepackage{algpseudocode}
\usecolortheme{default}

\renewcommand{\algorithmiccomment}[1]{\hfill\# \textit{#1}}
\algnotext{EndFor}
\algnotext{EndProcedure}
\algnotext{EndIf}
\algnotext{EndWhile}
\setbeamertemplate{footline}{%
  \hfill\insertframenumber/\inserttotalframenumber\hspace{0.5cm}\vspace{0.1cm}%
}
\setbeamertemplate{itemize item}[circle]
\setbeamertemplate{itemize subitem}[circle]
\setbeamercolor{itemize subitem}{fg=red}


\title{Jarníkův algoritmus}
\author{Jiří Prokop}
\institute{VUT FIT}
\date{2023}

\begin{document}

\frame{\titlepage}

\begin{frame}{Obsah}
    \tableofcontents
\end{frame}

\section{Motivace}
\begin{frame}
\frametitle{Problém}
\begin{itemize}
    \item Hledáme minimální kostru váženého neorientovaného grafu.
    \item minimální: nejmenší možná
    \item Co je kostra?
    \begin{quotation}
Kostrou grafu budeme rozumět libovolný podgraf, který hranami spojuje všechny vrcholy původního grafu a zároveň sám neobsahuje žádnou kružnici(jde o strom).
    \end{quotation}
    \item vážený: přechody mají cenu přechodu/váhu
    \item neorientovaný: přechody jde 'projít' oběma směry
  \end{itemize}
    \begin{figure}
        \includegraphics[width=0.4\textwidth]{Minimum_spanning_tree.svg.png}
        \caption{Minimalni kostra grafu}
    \end{figure}
\end{frame}

\section{Jarníkův algoritmus}
\begin{frame}{Jarníkův algoritmus}
    \begin{itemize}
        \item Ve světě znám jako \emph{Prim's alg.}
        \item Objeven Vojtěchem Jarníkem v roce 1930.
        \item Znovuobjeven v roce 1957 R. C. Primem a v roce 1959 E. W. Dijkstrou.
        \item Využívá tzv.'greedy' přístup - postupně přidává nejlevnější hrany.
        \item Složitost algoritmu závisí na použité datové struktuře.
    \end{itemize}
\end{frame}

\subsection{pseudokód}
\begin{frame}{Jarníkův algoritmus - pseudokód}
    \begin{figure}
        \begin{algorithmic}[0]
            \Procedure{prim}{graf}
            \State mst = prázdná množina 
            \Comment{pro držení vrcholů minimálním kostry} 

            \State startVertex = vrchol grafu
            \State mst.add(startVertex)
            \State hrany = hrany vrcholu startVertex
    
            \While{$mst.vertex.count < graf.vertex.count$}
                % \Comment{najdi minimální hranu v sadě hran}
                \State minEdge, minWeight = najítMinEdge(hrany)
                
                \State mst.add(minEdge)
                  
                \For{hrana in hrany spojené k minEdge}
                 % \Comment{přidej hrany připojené k vrcholu do sady hran}
                    \If{hrana is not in mst}
                        \State hrany.add(hrana)
                    \EndIf
                \EndFor
    
                \\ \Comment{Odstraňte minimální hranu ze sady hran, které chcete zvážit}
                \State hrany.remove(minEdge)
            \EndWhile
    
            \Return mst jako array
            \EndProcedure
        \end{algorithmic}
    \end{figure}
\end{frame}

\subsection{složitost}
\begin{frame}{Jarníkův algoritmus - složitost}
\framesubtitle{implementace pomocí haldy}
    \begin{itemize}
        \item časová složitost: \\
            $O(E \log V)$
            \begin{itemize}
                \item $E$ - počet hran 
                \item $V$ - počet vrcholů grafu
            \end{itemize}
        \item prostorová složitost: \\
            $O(V + E)$
            \begin{itemize}
                \item $O(V)$ pro držení minimální kostry
                \item $O(E)$ pro haldu
            \end{itemize}
    \end{itemize}
\end{frame}

\begin{frame}{Jarníkův algoritmus - složitost}
\framesubtitle{implementace pomocí tzv. Fibonacciho haldy}
    Je asymptoticky rychlejší, pokud je graf dostatečně hustý, takže platí, že $E$ je $\Omega(V)$.
    \begin{itemize}
        \item časová složitost: \\
            $O(E + V \log V)$
            \begin{itemize}
                \item $E$ - počet hran 
                \item $V$ - počet vrcholů grafu
            \end{itemize}
        \item prostorová složitost: \\
            $O(V + E)$ (stejná)
    \end{itemize}
\end{frame}

\subsection{algoritmus - příklad}
\begin{frame}{Jarníkův algoritmus - algoritmus}
Ukázka alg. na tomto grafu:
    \begin{figure}
        \centering
        \includegraphics[width=\textwidth]{Minimum-Spanning-tree-example-1.png}
    \end{figure}
\end{frame}

\begin{frame}{Jarníkův algoritmus - algoritmus}
První iterace - začínáme ve vrcholu 2
    \begin{figure}
        \centering
        \includegraphics[width=\textwidth]{Minimum-Spanning-tree-example-2.png}
    \end{figure}
\end{frame}

\begin{frame}{Jarníkův algoritmus - algoritmus}
Druhá a třetí iterace
    \begin{figure}
        \centering
        \includegraphics[width=0.75\textwidth]{Minimum-Spanning-tree-example-3.png}
    \end{figure}

    \begin{figure}
        \centering
        \includegraphics[width=0.75\textwidth]{Minimum-Spanning-tree-example-4.png}
    \end{figure}
\end{frame}

% \begin{frame}{Jarníkův algoritmus - algoritmus}
%  iterace
%     \begin{figure}
%         \centering
%         \includegraphics[width=0.5\textwidth]{Minimum-Spanning-tree-example-4.png}
%     \end{figure}

%     \begin{figure}
%         \centering
%         \includegraphics[width=0.5\textwidth]{Minimum-Spanning-tree-example-5.png}
%     \end{figure}

%     \begin{figure}
%         \centering
%         \includegraphics[width=0.5\textwidth]{Minimum-Spanning-tree-example-7.png}
%     \end{figure}
% \end{frame}

\begin{frame}{Jarníkův algoritmus - algoritmus}
    \begin{columns}
    \hspace{0.2cm}
        \begin{column}{0.5\textwidth}
            4. iterace
            \begin{figure}
                \includegraphics[width=\textwidth]{Minimum-Spanning-tree-example-4.png}
            \end{figure}
            5. iterace
            \begin{figure}
                \includegraphics[width=\textwidth]{Minimum-Spanning-tree-example-5.png}
            \end{figure}
        \end{column}
        \begin{column}{0.6\textwidth}
            \hspace{1.2cm} výsledná minimální kostra
            \begin{figure}
                \raggedleft
                \includegraphics[width=\textwidth]{Minimum-Spanning-tree-example-7.png}
            \end{figure}
        \end{column}
    \end{columns}
\end{frame}


\end{document}


